\section{Graded resolutions and leading coefficient spaces}
\label{sec:filtered}

We first prove a comparison for quotients of a free module whose monomial
positions have vector-space coefficients.  The application in the next
section will use a division algebra as that coefficient space.  The
comparison must retain each internal degree separately, so we work with
the graded Koszul complex throughout.

Let $F$ be a field and let $A=F[y_1,\ldots,y_r]$, with each variable of
degree one.  For a finitely generated graded $A$-module $B$, put
\[
 \beta^A_{p,j}(B)=\dim_F\Tor^A_p(B,F)_j,
 \qquad F=A/(y_1,\ldots,y_r).
\]
We use the shift convention $B(-s)_d=B_{d-s}$.  All modules considered
below are bounded below and have finite-dimensional graded components.

\begin{lemma}[Minimal graded resolutions]
\label{filt:minimal}
Every finitely generated graded $A$-module $B$ has a graded free resolution
whose differentials have entries in $A_{>0}$.  In such a resolution, the
number of free generators of degree $j$ in homological position $p$ is
$\beta^A_{p,j}(B)$.  If $H\subsetneq A$ is a homogeneous ideal, then
\begin{equation}
 \beta^A_{p+1,j}(A/H)=\beta^A_{p,j}(H)
 \qquad(p\geq 0,\ j\in\Z).
 \label{filt:ideal-quotient}
\end{equation}
If $H$ is monomial, a minimal resolution of $A/H$ can be chosen
homogeneous for the exponent multigrading, with all generator
multidegrees in $\Z_{\geq0}^r$.
\end{lemma}

\begin{proof}
Lift a homogeneous $F$-basis of $B/A_{>0}B$ to $B$ and map the free module
on these lifts to $B$.  This map is surjective: in each degree, an element
not yet expressed in the lifts differs from such an expression by a sum
of positive-degree multiples of elements of smaller degree, to which
induction applies.  Its kernel is finitely generated because $A$ is
Noetherian.  Repeating the construction for successive kernels gives a
free resolution.  At each stage, reduction modulo $A_{>0}$ identifies the
chosen free generators with a basis of the corresponding module modulo
$A_{>0}$, so the next kernel lies in $A_{>0}$ times that free module.
Thus every differential between free modules has positive-degree
entries.  Tensoring with $F$ makes these differentials zero and proves
the assertion about $\Tor$.

A minimal resolution of a proper homogeneous ideal $H$, followed by
the inclusion $H\hookrightarrow A$ and the quotient $A\twoheadrightarrow
A/H$, is a minimal resolution of $A/H$: the additional entries also
have positive degree since $H\subseteq A_{>0}$.  This proves
\eqref{filt:ideal-quotient}, including $H=0$.

For a monomial quotient, carry out the same construction with
multihomogeneous lifts.  The initial free module is $A$ in multidegree
zero.  A free module whose generator multidegrees are nonnegative has
no component in a multidegree outside $\Z_{\geq0}^r$; neither does its
kernel or a quotient of that kernel.  Induction therefore gives the
last assertion.
\end{proof}

\begin{lemma}[Koszul calculation]
\label{filt:koszul}
Let $K_\bullet(B)$ be the Koszul complex on $y_1,\ldots,y_r$ with
coefficients in $B$.  Its term of internal degree $j$ in position $p$ is
\[
 K_p(B)_j=\bigoplus_{\substack{\sigma\subseteq\{1,\ldots,r\}\\
                              |\sigma|=p}} B_{j-p}e_\sigma,
\]
where $e_\sigma$ has degree $p$.  Its homology satisfies
\[
 \dim_F H_p(K_\bullet(B)_j)=\beta^A_{p,j}(B).
\]
\end{lemma}

\begin{proof}
For $\sigma=\{i_1<\cdots<i_p\}$ the differential is
\[
 d(be_\sigma)=\sum_{q=1}^p(-1)^{q-1}y_{i_q}b\,
                    e_{\sigma\setminus\{i_q\}}.
\]
The complex $K_\bullet(A)$ is the tensor product over $F$ of the
one-variable complexes
\[
 0\longrightarrow F[y_i](-1)
   \xrightarrow{\ y_i\ }F[y_i]\longrightarrow0.
\]
Each has homology $F$ in position zero and no other homology.  A complex
of vector spaces splits as its homology with zero differential plus
contractible two-term summands.  Tensoring these splittings shows that
$K_\bullet(A)$ is a free resolution of $F$.

To compare this computation with a resolution of $B$, take any graded
free resolution $F_\bullet\to B$ and form the first-quadrant double
complex $F_\bullet\otimes_A K_\bullet(A)$.  Augmenting either factor
gives maps from its total complex to $F_\bullet\otimes_A F$ and to
$B\otimes_A K_\bullet(A)$, respectively.  The kernels have exact columns
or exact rows, since the terms of the other factor are free.  Their
total complexes are exact.  For exact columns, eliminate a cycle's
component of greatest horizontal index by a vertical boundary; the
correction introduces terms only in the next column to the left.
There are finitely many components on each total-degree diagonal,
so iteration eliminates the cycle.  The argument for exact rows
interchanges the two indices.  Both augmentation maps consequently induce
homology isomorphisms.  This identifies the homology of
$B\otimes_A K_\bullet(A)=K_\bullet(B)$ with $\Tor^A(B,F)$, preserving
the internal grading.
\end{proof}

We now order monomial positions without choosing an order on their
coefficients.  Let $V$ be a finite-dimensional $F$-vector space in
degree zero, and set
\[
 E=V\otimes_F A=\bigoplus_{\alpha\in\Z_{\geq0}^r}Vy^\alpha.
\]
Choose a total order $\prec$ on the exponents of each fixed total
degree, compatible with addition: if $\alpha\prec\beta$, then
$\alpha+\gamma\prec\beta+\gamma$ for every
$\gamma\in\Z_{\geq0}^r$.  The first nonzero position of an element is
its leading position in this order.

For a homogeneous submodule $W\subseteq E$ and an exponent $\alpha$,
define
\[
 V_\alpha=\left\{v\in V:
   vy^\alpha+\sum_{\substack{|\gamma|=|\alpha|\\\alpha\prec\gamma}}
        v_\gamma y^\gamma\in W
        \text{ for some }v_\gamma\in V\right\}.
\]
This is an $F$-subspace of $V$.  The resulting leading submodule is
\[
 \operatorname{in}_{\prec}W
       =\bigoplus_{\alpha\in\Z_{\geq0}^r}V_\alpha y^\alpha.
\]
Indeed, multiplication of a witnessing element by $y_i$ shows that
$V_\alpha\subseteq V_{\alpha+e_i}$, so this direct sum is an
$A$-submodule.

\begin{proposition}[Leading-submodule comparison]
\label{filt:comparison}
\label{filt:initial}
With $A,V,E,W$ and $\prec$ as above, for every $d\in\Z$, every
$p\geq0$, and every $j\in\Z$ one has
\begin{align}
 \dim_F(E/W)_d
    &=\dim_F(E/\operatorname{in}_{\prec}W)_d,
       \label{filt:hilbert}\\
 \beta^A_{p,j}(E/W)
    &\leq\beta^A_{p,j}(E/\operatorname{in}_{\prec}W).
       \label{filt:betti}
\end{align}
\end{proposition}

\begin{proof}
Write $Q=E/W$.  In degree $d\geq0$, let $E_{\succeq\alpha}$ be the
sum of the positions $Vy^\gamma$ with $|\gamma|=d$ and
$\alpha\preceq\gamma$, and define $E_{\succ\alpha}$ similarly with
strict inequality.  Their images in $Q_d$ give a finite descending
filtration.  Its quotient at $\alpha$ is
\begin{equation}
 \frac{\operatorname{im}(E_{\succeq\alpha}\to Q_d)}
      {\operatorname{im}(E_{\succ\alpha}\to Q_d)}
 \simeq
 \frac{E_{\succeq\alpha}}
      {E_{\succ\alpha}+(W_d\cap E_{\succeq\alpha})}
 \simeq (V/V_\alpha)y^\alpha.
 \label{filt:slot-quotient}
\end{equation}
The last isomorphism takes the coefficient at $\alpha$.  Summing its
dimensions proves \eqref{filt:hilbert}; both sides vanish for $d<0$.

Fix an internal degree $j\geq0$.  We give all terms of
$K_\bullet(Q)_j$ a common filtration indexed by the degree-$j$
exponents.  For a subset $\sigma$, put
$\mathbf{1}_\sigma=\sum_{i\in\sigma}e_i$.  In the summand with
exterior label $e_\sigma$, place the coefficient position $\alpha$ at
\[
                    \alpha+\mathbf{1}_\sigma.
\]
More explicitly, the filtration subspace at $\beta$ in that summand
is the image of the span of all $Vy^\alpha$ satisfying
$\beta\preceq\alpha+\mathbf{1}_\sigma$.  The full filtration subspace
is the direct sum over the labels $\sigma$.  Addition compatibility
means that the relative order of the positions within one summand
agrees with their order in $Q_{j-|\sigma|}$.

A differential deleting $i\in\sigma$ multiplies its coefficient by
$y_i$.  It sends the position with exponent $\alpha$ to that with
exponent $\alpha+e_i$, and
\[
 (\alpha+e_i)+\mathbf{1}_{\sigma\setminus\{i\}}
       =\alpha+\mathbf{1}_\sigma.
\]
Thus the differential preserves the filtration.  By
\eqref{filt:slot-quotient}, its map on the corresponding associated
graded positions is, up to the Koszul sign,
\[
 V/V_\alpha\longrightarrow V/V_{\alpha+e_i},
 \qquad [v]\longmapsto[v].
\]
This is well defined by $V_\alpha\subseteq V_{\alpha+e_i}$.
Consequently the associated graded complex is exactly
$K_\bullet(E/\operatorname{in}_{\prec}W)_j$, decomposed by the total
exponent $\beta$.

For completeness, if $d:C\to C'$ is a filtration-preserving linear
map between finite-dimensional filtered spaces, the filtration induced
on $\ker d$ gives an injection
\[
          \operatorname{gr}(\ker d)
               \hookrightarrow\ker(\operatorname{gr}d).
\]
Indeed a leading class represented by an element of $\ker d$ maps to
zero, and the induced filtration on the kernel makes this inclusion
injective.  Since taking associated graded preserves dimension, it
follows that
$\operatorname{rank}d\geq\operatorname{rank}(\operatorname{gr}d)$.
At a fixed term of a complex, the homology dimension is the term
dimension minus the ranks of its incoming and outgoing maps.  Applying
the rank inequality to these two maps proves that passing to the
associated graded complex cannot decrease this homology dimension.
Lemma~\ref{filt:koszul} now gives \eqref{filt:betti}.  For $j<0$ all
the relevant Koszul terms vanish, so the same assertion holds.
\end{proof}

Taking $V=F$ recovers the usual comparison with a monomial initial
ideal; see also \cite[Section~2]{MM}. Keeping $V$ arbitrary is essential
for the next application:
the coefficient algebra will enlarge all dimensions by the same
finite factor, which can then be cancelled.
